Nel corso degli ultimi anni, con l'avvento del Cloud Computing, si è assistito ad un importante cambiamento di paradigma nel mondo dell'Information Technology. Il Cloud Computing ha ridefinito, andando a semplificarne l'utilizzo, quelle che sono le operazioni che vengono effettuate per poter gestire i nostri sistemi informatici. \\
In questa prima parte del capitolo introduttivo è fornita una panoramica generale sulle caratteristiche e sull'importanza del Cloud Computing.

\section{Cloud Computing}
Non esiste un'unica definizione di Cloud Computing, infatti quest'ultimo viene definito in modo differente da varie aziende e sviluppatori. Una definizione data da Amazon, uno dei più importanti Cloud Provider del mondo, è la seguente: \\
\textit{"Il cloud computing consiste nella distribuzione on-demand delle risorse IT tramite Internet, con una tariffazione basata sul consumo. Piuttosto che acquistare, possedere e mantenere i data center e i server fisici, è possibile accedere a servizi tecnologici, quali capacità di calcolo, storage e database, sulla base delle proprie necessità affidandosi a un fornitore cloud quale Amazon Web Services (AWS)."} \cite{aws}. \\
Da questa definizione si evince che un Cloud è un modello di calcolo distribuito che permette di allocare, su richiesta (on-demand), un insieme di risorse di calcolo messe a disposizione da un fornitore, il Cloud Provider. \\
Il motivo principale per il quale questo paradigma si è diffuso molto rapidamente è il modello di costo su cui si basa il Cloud Computing, ovvero il modello pay-as-you-go. Con il Cloud le aziende riescono ad eliminare tutti quei costi che derivano dall'allocare una quantità maggiore (o minore) di risorse rispetto a quante siano effettivamente necessarie, infatti, grazie al modello pay-as-you-go, il mercato raggiunge uno stato in cui si incontrano domanda e offerta. Questo vantaggio non va solo a favore di quelle che sono le grandi imprese preesistenti, ma favorisce anche l'ingresso nel mercato di nuove imprese. \\ \\
Alla base del funzionamento del Cloud Computing vi è la virtualizzazione. La virtualizzazione consente ai Cloud Provider di offrire ai propri clienti risorse di elaborazione on-demand in una gamma molto ampia di configurazioni. Quando viene richiesta una macchina virtuale (VM) con delle particolari caratteristiche hardware (ad esempio velocità del processore, capacità di memoria...), il Cloud Provider non rilascia una macchina che ha esattemente quelle caratteristiche, piuttosto, grazie appunto alla virtualizzazione che consente di suddividere le risorse fisiche di una macchina in risorse virtuali più piccole, ricava le risorse necessarie da dispositivi esistenti più grandi \cite{practitioner}. \\ \\
Le architetture di Cloud Computing definiscono chi ha il controllo dell'infrastrurra e dove essa risiede. Le architetture possono essere suddivise in tre modelli:
\begin{description}
    \item[Public Cloud]Modello principale d'erogazione di servizi accessibili ad un pubblico generale su abbonamento, tramite Internet. I leader del settore sono Amazon Web Service, Microsoft Azure, Google Cloud. 
    \item[Private Cloud]Sistema di server dedicato completamente alle infrastrutture di una singola azienda. Possiede le stesse caratteristiche del Public Cloud ma offre una maggiore sicurezza.
    \item[Hybrid Cloud]Formato dalla combinazione di due o più infrastrutture Cloud distinte, in modo da poter ottenere i vantaggi delle varie infrastrutture.

\end{description}    
I vantaggi del Cloud Computing sono numerosi ed ognuno di essi avvalora il motivo per il quale il Cloud ha portato ad un importante cambiamento di paradigma nel mondo dell'Information Technology. Tra i più importanti vantaggi troviamo:
\begin{itemize}
    
    \item \textit{agilità} - il Cloud permette di poter avere piena libertà di sperimentare nuove idee grazie alla possibilità di poter accedere in modo semplice a diverse tecnologie in pochissimi minuti, usufrendo solo della quantità di risorse di cui si necessita;
    \item \textit{elasticità} - grazie al Cloud è sempre possibile effettuare un ridimensionamento delle risorse in uso in modo da poter rispondere ai picchi di carico e di poter rilasciare risorse inutilizzate;
    \item \textit{costi} - con il Cloud è possibile eliminare i costi relativi all'hardware e al software e tutti quei costi derivanti dalla loro gestione e configurazione, in favore di una spesa variabile nettamente inferiore;
    \item \textit{affidabilità} - il Cloud aumenta l'efficienza per quanto riguarda la gestione dei dati, più nel particolare la gestione dei backup di quest'ultimi, andando a ridurre i costi derivanti dal loro ripristino in caso di emergenza;
    \item \textit{sicurezza} - ogni Cloud Provider garantisce la sicurezza fisica delle macchine in cui risiedono dati e applicazioni,xx grazie all'integrazione di tecnologie e controlli che rafforzano la sicurezza di quest'ultimi. Tuttavia, i Cloud Provider non sono responsabili della sicurezza a livello applicativo, che rimane a carico del cliente.

\end{itemize}

\subsection{Modelli di servizio}
Possiamo dividere i vari servizi offerti dai Cloud Provider in diversi modelli di servizio, andando in questo modo a differenziare i servizi in base al livello di astrazione offerto all'utente. Esistono tre modelli principali di Cloud Computing che verrano di seguito descritti \cite{nist}. Nel capitolo 1.1.2 verrà descritto un quarto modello di servizio, ovvero il modello Function-as-a-Service (FaaS).
I tre modelli principali sono:
\begin{description}
    \item[Infrastracture-as-a-Service]Con il modello IaaS il consumatore ha a disposizione le infrastrutture di cui necessita sotto forma di computer (virtuali o su hardware dedicato) delle quali è possibile specificare le caratteristiche hardware e software. L'utente ha pieno controllo di gestione delle risorse che utilizza, ma non controlla i processi che consentono il funzionamento della macchina.
    
    \item[Platform-as-a-Service]Con il modello PaaS il consumatore non deve più gestire l'infrastruttura sottostante. La gestione dell'hardware e del sistema operativo passa al Cloud Provider in modo da permettere al consumatore di concentrarsi solo sulla gestione delle applicazioni. 

    \item[Software-as-a-Service]Con il modello SaaS il consumatore ha la possibilità di utilizzare software pronti all'uso, eseguiti e gestiti dal Cloud Provider, tramite l'utilizzo del browser. Un esempio comune di applicazione SaaS è Gmail: Gmail permette all'utente di inviare e ricevere e-mail tramite un qualsiasi browser, trascurando ogni aspetto dell'infrastruttura sottostante.
\end{description}

\subsection{Serverless Computing}
Il modello Function-as-a-Service (FaaS), conosciuto anche sotto il nome di Serverless Computing, è diverso dagli altri modelli di Cloud Computing introdotti nel capitolo precedente, in quanto, con questo modello, il Cloud Provider è responsabile di gestire sia l'infrastruttura Cloud sia la scalabilità delle app. Grazie a questo modello, le applicazioni vengono avviate solo quando necessario, evitando di lasciare attiva l'infrastruttura Cloud anche quando quest'ultima non è in uso. Le applicazioni saranno comunque eseguite su un server, ma con una differenza: tutte le operazioni che riguardano la gestione e la manutenzione dell'infrastruttura necessaria saranno a carico del Cloud Provider. In questo modo, gli sviluppatori potranno concentrarsi su quello che sarà il prodotto finale, ignorando tutte le operazioni di gestione e funzionamento dei server. \\
Il Serverless Computing si basa sul concetto di funzione, un blocco di codice indipendente che viene eseguito su Cloud in modo totalmente trasparente all'utente, in risposta ad uno specifico evento. Il Serverless Computing segue il modello di programmazione event-driven: ogni funzione viene eseguita in risposta ad un evento e può a sua volta emetterne un altro. In questo modo, una funzione serverless può emmetere un evento che scatenerà l'esecuzione di un'altra funzione serverless, e così via, andando in questo modo a creare delle catene di esecuzioni. Grazie a questo meccanismo è possibile partizionare la logica di business in più funzioni, in modo da poter favorire il riuso del codice.

\subsection{Multi-cloud}
Il multi-cloud è una metodologia di distribuzione di servizi con la quale è possibile accedere ai servizi offerti da diversi Cloud Provider in una singola architettura. Grazie a questa metodologia è possibile scegliere fra un infinito numero di combinazioni di servizi e risorse per poter raggiungere un obiettivo, dando luogo a diversi vantaggi:
\begin{itemize}
    
    \item \textit{costo-efficienza} - avendo un'ampia scelta di combinazioni fra risorse e servizi, è più probabile trovare la scelta più adatta, sia da un punto di vista economico che da un punto di vista di efficienza;
    \item \textit{riduzione della dipendenza} - data la possibilità di usufruire di servizi e risorse da diversi fornitori, viene eliminata la dipendenza da un singolo Cloud Provider;
    \item \textit{tolleranza agli errori} - in caso di malfunzionamento da parte di un Cloud Provider è possibile replicare l'intera infrastruttura su un altro Provider;
    \item \textit{distribuzione su più paesi} - possibilità di scegliere l'area geografica in cui i dati saranno memorizzati fisicamente, optando per paesi che più si adattano alle esigenze individuali sulla base di regolamenti e leggi interne.

\end{itemize}
L'utilizzo del multi-cloud, tuttavia, incrementa notevolmente la complessità del sistema, in quanto è necessario gestire un numero sempre maggiore di modelli diversi.

\section{Database}
Oggigiorno, lavorare con i dati è una necessità in qualsiasi contesto e applicazione reale. I database sono una soluzione pratica per poter archiviare e gestire informazioni (o dati).
In questa seconda parte del capitolo introduttivo è fornita un'introduzione generale sui database, soffermandosi sull'importanza della nascita dei database NoSQL.

\subsection{Definizione di database}
Un database è un insieme di informazioni strutturate, archiviate elettronicamente in un sistema informatico. In genere, i dati all'interno di un database vengono rappresentati tramite una serie di tabelle per poter garantire l'efficienza di elaborazione. I dati presenti all'interno di una tabella possono essere facilmenti visualizzati, gestiti, modificati, aggiornati e rimossi tramite l'utilizzo di query sui dati. \\
I primi sistemi di database nacquero intorno agli anni '60. Con gli anni, questi sistemi si sono evoluti sempre di più per poter soddisfare al meglio quelle che erano le richieste degli utenti. I sistemi di database più noti vennero introdotti intorno agli anni '80; tra questi è importante citare i database relazionali, ancora oggi utilizzati in moltissime applicazioni. Gli elementi in un database relazionale sono rappresentati sotto forma di tabelle composte da colonne e righe. Questo tipo di rappresentazione offre la soluzione più efficiente e flessibile per poter gestire informazioni strutturate \cite{oracle}. \\
Più di recente sono comparsi i database NoSQL (non relazionale), in risposta all'aumento della mole di dati e alla velocità con cui vengono prodotti. A differenza dei database relazionali, i database NoSQL non utilizzano un singolo modo per poter rappresentare i dati, ma offrono una varietà di modelli con i quali è possibile rappresentarli e gesterli.

\subsection{Database NoSQL}
Alcune persone affermano che la parola NoSQL stia per "non SQL", altre invece per "non solo SQL".  In entrambi i casi, con il termine "database NoSQL", si fa riferimento a un qualsiasi database non relazionale, ovvero un modello di database dove i dati vengono rappresentati in un formato diverso dalle tabelle. Un errore molto comune ricade nell'interpretare in modo errato il significato del termine "non relazionale": è sempre possibile mettere in relazione fra loro i dati presenti in un database non relazionale, utilizzando però una strategia diversa \cite{mongodb}.\\
I motivi principali che hanno portato ad un utilizzo sempre più frequente dei database NoSQL sono:
\begin{itemize}
    
    \item \textit{flessibilità} - i database NoSQL offrono agli sviluppatori la possibilità di poter rappresentare i dati utilizzando diversi modelli di database;
    \item \textit{elevate prestazioni} - i database NoSQL sono ottimizzati per modelli di dati specifici e schemi di accesso che consentono prestazioni più elevate rispetto ai database relazionali;
    \item \textit{distribuzione} - la flessibilità nella clusterizzazione e nella replicazione dei dati permette di distribuire su più nodi lo storage;
    \item \textit{funzionalità} - i dati vengono archiviati e richiamati attraverso API basate su oggetti.

\end{itemize}
\subsection{Differenza fra NoSQL e SQL}
La differenza principale fra i database NoSQL e i database SQL è rappresentata dal modello di scalabilità adottato. I database NoSQL sono progettati per la \textbf{scalabilità orizzontale}. Scalare orizzontalmente (scale-out) significa ovviare alla saturazione di un nodo di un sistema aggiungendone un altro dello stesso genere, ad esempio un altro server o un altro storage. Questo modello permette di ampliare la capacità e le performance di un sistema in modo teoricamente infinito, data la sempre possibilità di aggiungere un nuovo nodo all'interno di esso \cite{wisefrogs}. La scalabilità orizzontale trova un forte riscontro con quelli che sono i servizi offerti dal Cloud. Grazie al Cloud è sempre possibile aggiungere (o rimuovere), in modo del tutto automatico, nuovi nodi all'interno del sistema in risposta ai picchi di carico (o ad una riduzione di domanda). I database SQL sono progettati invece per la \textbf{scalabilità verticale} (scale-up). Essa si ottiene aumentando le risorse di un singolo nodo del sistema, ad esempio aggiungendo CPU ad un server o dischi ad uno storage. Lo svantaggio di questo modello consiste nel costo: l'aggiornamento hardware di una macchina può essere economicamente molto più gravoso dell'acquisto di un'ulteriore macchina di pari potenza. \\
Per consentire la scalabilità orizzontale, la maggior parte dei database NoSQL, non supporta le transazioni ACID (Atomicity, Consistency, Isolation, Durability):
\begin{itemize}
    \item \textit{atomicity} - le transazioni devono essere completate o fallite nel loro insieme;
    \item \textit{consistency} - per tutta la durata delle transazioni devono essere rispettati i \textbf{vincoli di integrità};
    \item \textit{isolation} - il fallimento di una transazione non deve interferire con le altre transazioni in esecuzione;
    \item \textit{durability} - i risultati di una transazione completata con successo sono resi permanenti nel sistema.
\end{itemize}
Esse vengono sostituite con il modello BASE (Basically Available, Soft State, Eventually Consistant):
\begin{itemize}
    \item \textit{basically available} - il sistema deve garantire la disponibilità dei dati;
    \item \textit{soft state} - la consistenza dei dati è un problema dello sviluppatore e non deve essere gestita dal database;
    \item \textit{eventually consistant} - i risultati di una transazione completata con successo vengono propaghati gradatamente, un nodo alla volta, fino a far diventare consistente l'intero sistema.
\end{itemize}
Il modello BASE abbandona quindi i requisiti di Atomicity, Consistency e Durability per ottenere una migliore scalabilità. Vale quindi il principio secondo cui la disponibilità è prioritaria rispetto alla coerenza.

\subsection{Tipi di database NoSQL}
Come già introdotto nei capitoli precedenti, esistono diversi modi per rappresentari i dati all'interno di un database NoSQL. Nel corso degli anni, sono emerse tre metodologie principali:
\begin{description}
    \item[Database chiave-valore]I database chiave-valore immagazzinano i dati come un insieme di coppie di chiave-valore dove una chiave rappresenta un identificatore univoco. Questi database sono particolarmente adatti per casi d'uso in cui è necessario effettuare ricerche rapide di singoli blocchi di informazione.
    \item[Database documentali]I database documentali nascono con l'obiettivo di semplificare agli sviluppatori la ricerca e la memorizzazione di dati in un database. I dati all'interno di un database documentale vengono rappresentati tramite oggetti JSON, in modo da permettere agli sviluppatori di mantenere lo stesso modello di rappresentazione dei dati sia nel codice dell'applicazione sia all'interno del database. L'utilizzo di questa metodologia permette di ottenere un notevole aumento delle presentazioni nell'esecuzione di query sul database, a discapito di un aumento della ridondanza. Di seguito è riportato un esempio per comprendere al meglio quanto appena affermato. \\
Supponiamo di voler memorizzare all'interno del nostro database un insieme di libri. In un database relazionale il record di un libro viene normalizzato e archiviato in tabelle distinte; in questo contesto il record di un libro viene diviso in due tabelle: una relativa al libro e l'altra relativa all'autore del libro. L'operazione di normalizzazione è un procedimento volto all'eliminazione della ridondanza informativa, permessa dai vincoli di integrità referenziale, con la quale è possibile creare delle relazioni fra le tabelle. I database NoSQL di tipo documentale, molto spesso, non supportano vincoli di integrità referenziale, e quindi non permettono la creazione di relazioni fra i dati Questo implica che le informazioni riguardanti il libro e l'autore devono necessariamente essere inserite in un singolo oggetto. \\
Consideriamo il caso in cui volessimo memorizzare un nuovo libro scritto da un autore già presente all'interno del database. Nel database relazionale ci basterà aggiungere un record all'interno della tabella dei libri, andando a creare una relazione con il record appena aggiunto e il record relativo all'autore del libro già presente all'interno della tabella relativa alle informazioni degli autori. All'interno del database NoSQL di tipo documentale è necessario creare un nuovo oggetto che conterrà le informazioni sia del libro e sia dell'autore del libro. In un database NoSQL, quindi, risulta necessario memorizzare dati ridondati all'interno del database. Anche se da una parte la ridondanza dei dati produce mancanza di efficienza nella gestione e nell'occupazione della memoria, in questo modo è possibile ottenere tutte le informazioni riguardanti un libro semplicemente accedendo ad un singolo oggetto. L'utilizzo di un database relazionale comporta la necessità di accedere a due tabelle differenti, e quindi un tempo di esecuzione più lento rispetto all'utilizzo di un database documentale. Ovviamente, in contesti reali, i dati vengono normalizzati in molte più tabelle, comportando quindi un notevole incremento del tempo di esecuzione delle query.\\
Nel capitolo successivo verrà introdotto MongoDB, un database distribuito di tipo documentale. Il suo utilizzo risulterà chiave nello sviluppo di questa tesi.
    \item[Database a grafo]I database a grafo nascono con l'obiettivo di facilitare la creazione e l'esecuzione di app che lavorano con set di dati ad elevata connesione.  All'interno del database i dati vengono rappresentati tramite i vertici di un grafo e le varie connessione tramite degli archi. Grazie a questo modello di rappresentazione, a differenza del modello documentale e quello relazionale, è possibile avere un idea chiara e intuitiva delle varie relazioni presenti all'interno del database.
\end{description}

\subsection{MongoDB}
MongoDB è un database NoSQL di tipo documentale open source che fornisce supporto per sistemi di storage destinato agli sviluppatori di applicazioni moderne e al cloud. MongoDB utilizza un meccanismo di query JSON facilmente componibili che consente di filtrare e ordinare i dati in base a qualsiasi campo, indipendentemente dal livello di annidamento all'interno del documento. \\
MongoDB, diversamente da molti altri database di tipo documentale, supporta le transazioni ACID complete e permette due tipi di realazioni anzinché una: riferimento e incorportamento. \\
MongoDB negli anni è diventato uno dei database più diffusi per diverse ragioni:
\begin{itemize}
    
    \item \textit{completamente automatizzato} - il servizio di database offerto da MongoDB è completamente gestito da MongoDB Atlas. Il provisioning, la configurazione e la distribuzione dell'infrastruttura sono completamente automatizzati con Atlas. Grazie a queste funzionalità, Atlas permette di utilizzare in modo semplice e intuitivo un database NoSQL su Cloud, in modo da concentrarsi solo su quello che è la memorizzazione e la gestione dei dati, ignorando tutta la gestione dell'infrastruttura sottostante. Lo sviluppatore Ken Hoff afferma che: \textit{"Il bello della gestione di MongoDB con MongoDB Atlas è che praticamente quasi tutto il tempo, non dobbiamo preoccuparcene. Non solo possiamo apportare rapidamente modifiche alla nostra applicazione sfruttando il modello dati flessibile di MongoDB, ma possiamo anche implementare qualsiasi modifica del database downstream in tempo reale o far girare facilmente nuovi cluster per testare nuove idee. Tutto ciò può accadere senza influire sulle attività di produzione; non ci sarà da preoccuparsi del provisioning dell’infrastruttura, della configurazione dei backup, del monitoraggio, ecc. È una vera bellezza."} \cite{mongoefficiency};
    \item \textit{compatibile} - MongoDB è completamente compatibile con i vari servizi di gestione di database documentali offerti dai vari Cloud Provider;
    \item \textit{multi-cloud} - MongoDB è l'unico database multi-cloud distribuito a livello globale. Permette la possibilità di sviluppare un singolo cluster su diversi Cloud Provider senza la complessità operativa della gestione della replica e della migrazione dei dati tra cloud;
    \item \textit{affidabile} - MongoDB è stato creato con tolleranza di errore distribuita e ripristino automatizzato dei dati. I cluster sono progettati per la resilienza con un set di repliche a tre nodi minimo distribuiti tra le zone di disponibilità all'interno di un'area cloud. Nel caso in cui un nodo primario non è disponibile, Atlas eseguirà automaticamente un failover in pochi secondi, senza necessità di intervento manuale.

\end{itemize}